\documentclass[runningheads,a4paper]{llncs}

\usepackage[T1]{fontenc}
\usepackage{lmodern}
\usepackage{fullpage}
\usepackage{microtype}
\usepackage{amsmath,amssymb,mathtools}
\usepackage[hidelinks]{hyperref}
\usepackage[nameinlink,noabbrev]{cleveref}

% Shared notation belongs here once it has been fixed and reviewed.
\newcommand{\bits}{\{0,1\}}
\newcommand{\negl}{\mathop{\mathrm{negl}}}
\newcommand{\poly}{\mathop{\mathrm{poly}}}
\newcommand{\GC}{\mathsf{GC}}
\newcommand{\sample}{\mathrel{\leftarrow}}


\title{The YASO Model: YOSO without Erasure}
\titlerunning{The YAMSO Model}
\author{Jesper Buus Nielsen\inst{1} \and Ivan Damg{\aa}rd\inst{1} \and
  Sannidhi V. Hebbar\inst{2}}
\authorrunning{J. B. Nielsen et al.}
\institute{Aarhus University, Denmark\\
  \email{\{jbn,ivan\}@cs.au.dk}
  \and
  Indian Institute of Science, India\\
  \email{sannidhiv@iisc.ac.in}}

\begin{document}

\maketitle

\begin{abstract}
The YOSO paradigm protects against a strongly adaptive adversary. Each 
ephemeral role prepares one message, erases all remaining state, and only
then speaks. We ask what remains possible when erasure is unavailable and
speaking instead exposes the sender's complete state. At first sight, such a
model appears useless. Nevertheless, it has a different irreversible
resource: a machine may choose never to speak, in which case the secret state
assigned to it can remain hidden. We call the resulting model YAMSO, for ``You
At Most Speak Once.'' This preliminary article introduces the model, gives a
finite two-phase garbled-circuit blueprint, and develops a candidate
prefix-preserving garbled-circuit factory for ongoing computation. The latter
uses one-use common-input certificates to bind a copyable obfuscated stream to
the canonical ledger history. We isolate its security into three explicit
replacement lemmas. A horizon-free public-program realization and online
random-oracle programming are the two main remaining gaps.
\keywords{adaptive security \and YOSO \and erasure-free protocols \and
one-time roles \and garbled circuits \and indistinguishability obfuscation}
\end{abstract}

\section{Introduction}
\label{sec:introduction}

The YOSO (You Only Speak Once) model was introduced as a way to obtain secure
multiparty computation against a very strong adaptive adversary
\cite{GentryEtAl2021YOSOFull}. A protocol role is assigned to an ephemeral
party which prepares one message, deletes the secret state used to prepare
it, and then sends the message. Only at this point does the role reveal
itself. Corrupting the sender after publication is therefore too late: the
message is already available and the secrets that produced it no longer
exist. Later work has developed this anonymous, ephemeral-role paradigm in
several directions; the erase-before-send discipline remains central.

This paper starts from the deliberately unfriendly question: what if deletion
is impossible? The restriction might be legal---a regime may require state
to be retained for later inspection---or technological. For example, a
party running in a cloud environment might fear that speaking identifies its
machine and enables the cloud operator to inspect a residual snapshot despite
an attempted deletion. We capture the extreme version of this concern by
declaring that whenever an ordinary party speaks, its complete state leaks to
the adversary.

The immediate answer seems to be ``obviously impossible.'' Eventually all
parties relevant to the computation will speak and hence reveal everything.
This conclusion overlooks one asymmetry. Speaking is optional. A machine may
remain silent forever, in which case its secret state may stay hidden. We
therefore place mutually exclusive pieces of state behind different
public-key roles. The selected roles speak and are expended; the unselected
roles never speak. In this sense, silence replaces erasure as the one-time
resource.

We call the resulting setting \emph{YAMSO}: You At Most Speak Once. In a
YAMSO execution, a synchronous ledger maintains a growing ordered list of
public keys and messages. The public-key index determines the protocol role.
For the preliminary model, each newly registered computation role is corrupted
independently with probability $p$, whereas named application parties retain
the ordinary corruption interface. An honest role keeps its secret key inside
a protected ledger interface until its scheduled opportunity to speak. Once a
computation role speaks, its key and complete role state may become public;
its message is already irrevocable and the role is never used again.

The intended real-world interpretation has a much larger population of
apparently interchangeable machines. The adversary compromises a fixed
fraction of that population, while fresh clients and role states are placed
on machines randomly and without exposing the placement before a machine is
corrupted or speaks. Random placement then turns the fixed compromised set
into an approximately random intersection with each protocol committee. The
Bernoulli experiment in this paper deliberately collapses that mechanism into
a simple model; it does not model mobile corruption.

\paragraph{Finite feasibility.}
The basic two-phase blueprint isolates why the model need not be vacuous. For
each garbled input wire, the two structural labels are placed behind disjoint
committees, their semantic order is hidden, and an activation token causes
only one committee to speak. Its complete exposed state contains shares of a
label that is becoming public anyway. Honest members of the sibling committee
remain silent forever and protect the unused label. A missing or malformed
input closes onto a fixed structural branch whose hidden semantic value is a
uniform default bit.

\paragraph{Ongoing computation.}
The difficult part is not evaluating another garbled circuit; it is preserving
the hidden correlation between the current state-label pair and every later
stretch without giving a speaking linker both labels. We propose a finite
public \emph{oblivious garbled-circuit sampler} (OGCS), implemented
conjecturally by an obfuscated streaming program. It fixes one public
transition circuit, receives a growing sequence of certified public inputs,
and emits exactly prefix-preserving garbled stretches. The output-state label
pairs of one stretch are literally the input-state pairs of the next.

A copyable obfuscated program can otherwise be evaluated on invented public
keys and forked ledger histories. Our second use of silence is therefore a
one-time common-input certificate. Each bit of a high-distance encoding is
represented by two one-way preimages placed in different one-use roles. The
role matching the canonical message speaks; its opposite remains silent.
After the correct certificate is public, the speakers may reveal everything
without enabling a conflicting certificate. This observation removes the
long-lived certification key we initially appeared to need. The ledger still
provides a common immutable view, but it need not hold a signing secret.

This is an instance of the model changing the design intuition. If a single
certifier speaks, its leaked signing key destroys future authenticity. The
natural reaction is to declare certification impossible. YAMSO instead asks
which secret holders need never speak. The resulting construction is closely
related to known high-distance Lamport signatures, but its use as persistent
common truth under total post-speech exposure is the mechanism we need here.

\paragraph{Status and proof strategy.}
The ongoing construction is intentionally presented as a candidate, not a
theorem. We modularize its analysis by three whole-process replacement lemmas:
canonical one-use certification, programmable prefix-preserving OGCS, and
two-stage finite garbling/readout with full speaker exposure. This avoids a
naive UC-hybrid decomposition whose interfaces would export internal keys to
the environment. The first version fixes corrupt input providers before setup
and has public outputs. Two central realization gaps remain: the cited
general-purpose ordinary-iO-based succinct TM-iO is bounded-input rather
than horizon free, and predictable random-oracle points do not yet support
one online simulator. A hidden nonce prepared one
stretch ahead is the intended lift for the second gap.

\paragraph{Relation to prior work.}
Our starting point is YOSO \cite{GentryEtAl2021YOSOFull}. Anonymous blockchain
committees tolerating a corruption rate around one quarter are known, but
their secret handoff relies on erasures
\cite{BenhamoudaEtAl2020BlockchainSecret}. Blockchain one-time programs
already demonstrate conditional release of one garbled input-label vector on
an immutable ledger \cite{GoyalGoyal2017Blockchains}. No-erasure hidden
committees also appear in communication-local MPC, although there the selected
committee runs an ordinary multi-round computation
\cite{ChandranEtAl2024CommunicationLocalMPC}.

For certification, Katz and Vaikuntanathan give a direct high-distance
Lamport precedent \cite{KatzVaikuntanathan2009LeakageSignatures}, and Kelsey,
Lang, and Lucks study distributed hash-based signing and the
disjoint-honest-trustees one-time-key-reuse problem in the absence of
broadcast \cite{KelseyLangLucks2025DistributedHBS}. Neither result states
our one-shot common-input property under post-speech full exposure. For the
streaming component, relevant technical precedents include universal samplers,
cryptographic iterators, succinct and reusable garbling, TM obfuscation, and
streaming functional encryption
\cite{HofheinzEtAl2014UniversalSamplers,KoppulaLewkoWaters2014IOTM,%
LinPass2014SuccinctGarbling,GoldwasserEtAl2013TuringMachines,%
GuanKorbSahai2022StreamingFE}. We use these as component precedents rather
than claiming that an existing theorem already gives the YAMSO construction.

\section{Preliminaries: What Must Be Fixed}
\label{sec:preliminaries}

This section is presently a roadmap.  The final paper should reproduce every
definition whose details affect the YAMSO claim, while importing routine
cryptographic syntax and standard security definitions by reference.

\subsection{Conventions}

We write $\lambda$ for the security parameter, $[n]=\{1,\ldots,n\}$, and
$\bits^*$ for the set of finite bit strings.  All adversaries and protocols
are probabilistic polynomial time unless explicitly stated otherwise, and
$\negl(\lambda)$ denotes a negligible function.  Protocol identifiers,
epochs, wire identifiers, public-key indices, ledger positions, and object
types must be included in every cryptographic domain separator.  A computation
key is assigned to exactly one canonical protocol role.

\subsection{Definitions that belong in this paper}

The following objects are model-specific and should eventually be defined in
full.

\begin{description}
  \item[Ledger.] One synchronous, append-only list $L$ seen identically by
    every party.  It begins with a generic CRS, grows by registration and
    scheduled message entries, and performs no protocol-specific validation.
    It has no cryptographic certification key.

  \item[Growing key registry.] An adversarial registration request creates a
    computation PID equal to its fresh public key.  A registration from an
    existing application PID creates and publicly binds a protected key to
    that named party.  Neither request contains a protocol role; registration
    order determines roles at the protocol level.

  \item[YAMSO turn.] A public initialization parameter
    $\mathsf{WhosTurn}$ selects one unspent PID from the current $L$.  The
    selected party either produces one message or receives an explicit
    $\mathsf{NoMsg}$ entry and is then spent forever.

  \item[Protected local evaluation.] An honest scheduled party supplies an
    algorithm $A$ privately to the ledger, which evaluates $A(sk,L)$ without
    first returning $sk$ to the party.  This lets an ordinary role reveal its
    entire key and an input role reveal only an input-dependent function of
    its state through the same generic interface.

  \item[Two corruption rules.] Each computation PID is independently corrupt
    with probability $p$ when registered and, if honest, cannot later be
    corrupted.  Named application parties have the ordinary adaptive
    corruption interface in the model.  The first construction statement
    restricts their corrupt set to be fixed before setup.

  \item[Closed input slot.] At a protocol-defined cutoff, the public ledger
    prefix determines one immutable input record.  Public branch-indexed
    one-way tags determine whether it contains exactly one valid activation
    token.  A missing, ambiguous, or malformed record selects the structural
    default branch; late messages are ignored.

  \item[Common-input certificate.] Fresh one-use roles hold opposite one-way
    preimages for the coordinates of an encoded digest.  Honest roles compute
    the value to be certified from the same closed ledger cut, rather than
    signing a caller-supplied string.  Verification permits enough missing
    openings for liveness but must rule out a second accepted digest.

  \item[OGCS prefix.] A fixed finite public object is evaluated on complete
    finite histories.  Outputs on every accepting history are literal
    prefixes.  A history that first differs at some accepted command must not
    reveal the hidden logical state of the canonical history; a rejected
    history never extends again.
\end{description}

\subsection{Primitives that can largely be imported}

Standard definitions of pseudorandom functions, collision-resistant and
one-way hashing, commitments, secret sharing, public-key encryption, and
circuit garbling need not be repeated in full.  We should state the exact
variants used and cite them.  The following interfaces require particular
care.

\begin{description}
  \item[Role-state encryption.] Role membership is public in the preliminary
    model, so ordinary multi-user encryption can protect the state placed
    under each registered key; recipient anonymity is not required by this
    abstraction.  The semantic association between the two structural
    committees for an input wire must nevertheless remain hidden.  Each key
    has one canonical role packet.  Off-path evaluations may address further
    packets to the same key, and later disclosure decrypts them; the OGCS
    isolation lemma must make all such plaintexts independently rooted fork
    material.

  \item[Verifiable threshold sharing.] The construction needs privacy below
    the reconstruction threshold, reconstruction from honest shares, and
    public rejection of malformed shares or keys.  Public verification data
    must not disclose the semantic branch permutation.

  \item[Two-stage garbled circuits.] Preparation fixes external-input label
    pairs, tokens, and role packets before inputs arrive.  Completion occurs
    only after input closure and must support prescribed active external and
    carried-state labels, prescribed public output, hidden next-state label
    pairs, and consistency with every later-exposed selected role.  This
    causal interface is stronger than a bare one-shot privacy statement.

  \item[High-distance one-time opening.] The certificate uses the distance
    idea from Lamport-style leakage-resilient signatures
    \cite{KatzVaikuntanathan2009LeakageSignatures}.  The paper must prove its
    own partial-opening liveness and non-equivocation statement under random
    role corruption and complete post-speech exposure, including an explicit
    code-rate-versus-tail calculation.

  \item[Succinct TM-iO and puncturing.] For a declared input-length bound, the
    candidate OGCS can use a compact obfuscated Turing-machine program with
    puncturable PRF/KDF state
    \cite{KoppulaLewkoWaters2014IOTM,AnanthJainSahai2017TMIO}.  We must state
    the bound, succinctness, functionality-equivalence condition, and any
    equal-running-time requirement.  The cited general-purpose
    ordinary-iO-based bounded-input results do not by themselves realize the
    horizon-free replay interface of \cref{sec:ongoing}.

  \item[Delayed backdoor programming.] Hofheinz et al.'s adaptive universal
    sampler provides the hidden-trigger paradigm
    \cite{HofheinzEtAl2014UniversalSamplers}.  Selectively secure pseudorandom
    puncturable deterministic encryption gives a related formulation
    \cite{KoppulaPoelstraWaters2017FastSamplers}.  Their theorems do not apply
    automatically to our stateful syntax; the OGCS replacement lemma must be
    proved separately.

  \item[Statistical binding for iO hybrids.] Computational unforgeability does
    not by itself justify iO between programs that differ on an invalid but
    mathematically possible input.  In its binding setup, a somewhere
    statistically unforgeable signature makes alternate valid signatures at a
    selected round and message statistically nonexistent, except with
    negligible setup probability \cite{FernandoEtAl2023SSU}.  Adapting that
    property to our certificate chain is an open proof obligation.
\end{description}

\paragraph{Proof boundary.}
The main text uses the generic ledger interface of \cref{sec:model}.  Once the
API is stable, a detailed $\mathcal F_{\mathsf{Ledger}}$ can be placed in an
appendix.  The target functionality contains only named application parties;
the randomly corrupted computation population is implementation machinery to
be simulated away.  Intermediate cryptographic components are handled by
whole-process replacement lemmas, avoiding hybrid interfaces that export
their secret trapdoors to the environment.

\section{The YAMSO Ledger Model}
\label{sec:model}

We describe the ledger interface at the level needed for the finite
construction in \cref{sec:basic-construction}. Our approach follows the
modelling style of \cite[Section~2.1]{DinsdaleYoungEtAl2020Afgjort}: we state
an intentionally strong and simple abstraction, isolate the properties used
by the protocol, and postpone its detailed realization. A later version will
express this interface as an ideal functionality
$\mathcal F_{\mathsf{Ledger}}$, in the style of the total-order-broadcast
functionality of \cite{DamgardEtAl2025AsyncYOSO}. There is no reason to fix
that UC bookkeeping before the API itself has stabilized.

The ledger is generic. It knows nothing about universal sampling, garbled
circuits, committees, input wires, or the meaning of protocol roles. Those
belong to the protocol using the ledger.

\subsection{Public state and initialization}

Fix a security parameter $\lambda$, a session identifier $\mathsf{sid}$, a
constant $p<1/3$, a generic setup distribution $\mathsf{Setup}$, a
key-generation algorithm $\mathsf{KeyGen}$, and a description of a deterministic
protocol-dependent algorithm $\mathsf{WhosTurn}$. The description of
$\mathsf{WhosTurn}$ is supplied to the ledger as a fixed initialization value;
it cannot be changed in response to the execution. The ledger samples
\[
  \mathsf{crs}\leftarrow\mathsf{Setup}(1^\lambda)
\]
and initializes the public append-only list as
$L=((\mathsf{Setup},\mathsf{crs}))$. The concrete protocol gives meaning to
$\mathsf{crs}$; the ledger does not interpret it.

We assume synchronous computation and an unrealistically strong broadcast
abstraction. After an entry is appended, every party sees exactly the same
list $L$ at the same time. An append is atomic, cannot be suppressed after
it occurs, and fixes the position of that entry. There are no delayed,
private, or party-local ledger views. This deliberately idealized assumption
removes consensus and network scheduling from the present question.

Apart from its own setup, registration, and timeout records, the ledger does
not enforce protocol syntax. In particular, an arbitrary message supplied by
a corrupt scheduled party is appended verbatim. The protocol reading $L$
decides which entries are well-formed and ignores malformed ones.

\subsection{Two registration paths}

The registry grows during the execution. There are two ways to add a key;
neither command takes a protocol role as input. The interface imposes no
a-priori bound on the registry, although any efficient execution creates only
polynomially many keys and every finite protocol instance uses a finite
prefix.

\paragraph{Computation-party registration.}
On a $\mathsf{Register}$ request from the adversary, the ledger samples
\[
  (pk_i,sk_i)\leftarrow\mathsf{KeyGen}(1^\lambda),
\]
where $i$ is the next computation-key index, creates an ITM whose PID is
$pk_i$, and appends $(\mathsf{Register},pk_i)$ to $L$. Independently it
samples $c_i\leftarrow\mathsf{Bernoulli}(p)$. If $c_i=1$, the new
computation party is corrupt and $sk_i$ is given to the adversary. Otherwise
it is honest, the ledger retains $sk_i$ as protected state, and the adversary
cannot corrupt this party later. Registration is indivisible: after
requesting it, the adversary cannot discard the public key because it
dislikes the corruption outcome.

\paragraph{Application-party registration.}
A fixed application party with ordinary UC identity $P$ may also input
$\mathsf{Register}$. The ledger samples
$(pk_P,sk_P)\leftarrow\mathsf{KeyGen}(1^\lambda)$, stores the association
$P\leftrightarrow(pk_P,sk_P)$, and appends
$(\mathsf{Register},P,pk_P)$ to $L$. It does not return $sk_P$ to an honest
$P$. There is no Bernoulli corruption experiment for such a party: the
adversary may corrupt $P$ adaptively in the ordinary UC sense. If $P$ is
already corrupt, or becomes corrupt later, the ledger releases $sk_P$ and
its retained record of $P$'s protected calls to the adversary. The target
ideal functionality supplies the simulator with any input or output leakage
prescribed upon that corruption.

Each PID registers at most once per session. Application parties are the
meaningful principals at the interface of the computation---for example, an
account owner supplying an input. Their identity need not and generally
cannot be hidden. Computation parties, by contrast, are protocol machinery.

A public protocol-level rule maps registration positions to roles. Thus the
seventh computation key may represent the eighth share-release role. This
mapping is not an argument to $\mathsf{Register}$ and is not interpreted by
the ledger.

\subsection{Turns and protected local evaluation}

At each synchronous step, the fixed algorithm $\mathsf{WhosTurn}(L)$ returns
either a registered public key that has not previously had a turn, or
$\bot$. It may schedule a computation key $pk_i$ or the registered key
$pk_P$ of an application party. A well-formed protocol chooses
$\mathsf{WhosTurn}$ so that it waits for enough registrations and schedules
each required role in the appropriate order. The ledger does not interpret
why a particular key is scheduled.

When a key $pk$ is scheduled, its owner has one synchronous opportunity to
speak. Write $L_{\mathsf{pre}}$ for the list at the start of that turn. An
honest owner privately gives the ledger a finite polynomial-time algorithm
$A$. The ledger samples and retains any coins required by $A$, and evaluates
\[
  m\leftarrow A(sk,L_{\mathsf{pre}}),
\]
using the protected secret key associated with $pk$. If $m\neq\bot$, it
signs $(\mathsf{sid},L_{\mathsf{pre}},m)$ using the signing component of $sk$
and appends
\[
  (\mathsf{Message},pk,m,\sigma)
\]
to $L$. Neither $A$ nor the unrevealed part of $sk$ is output. The ledger
retains the algorithm and its coins as the protected call record rather than
modelling an erasure.

If the owner is corrupt, the adversary knows $sk$ and may supply any pair of
bit strings $(m,\sigma)$ during the turn. The ledger appends
$(\mathsf{Message},pk,m,\sigma)$ verbatim without checking a signature, a key
relation, or any other protocol-specific predicate. If an
honest $A$ returns $\bot$, or if no message is supplied by the end of the
turn, the ledger appends $(\mathsf{NoMsg},pk)$. This timeout is immediately
visible to everyone and lets $\mathsf{WhosTurn}$ advance past a silent role.
After either record is appended, $pk$ is spent and is never scheduled again.

For an honest nonempty message, protected evaluation, signing, and append are
one atomic operation. The adversary sees the result only after it is
irrevocably in $L$. The ledger does not give an adversarial scheduler an
opportunity to inspect the result and then suppress it.

We assume that $\mathsf{KeyGen}$ may provide a public relation
$\mathsf{KeyMatch}(pk,sk)$. A protocol that uses disclosure of $sk$ as its
message can then check that the disclosed key belongs to the scheduled
public key. If the chosen key system has no such relation, the ledger does
not manufacture one: the protocol must tolerate the missing check or choose
different key syntax.

\subsection{How protocols use the interface}

The following examples explain the intended use without adding commands to
the ledger.

\paragraph{Ordinary computation role.}
From $\mathsf{crs}$ and the public list, all parties may derive the same
protocol object, including a ciphertext $C_{pk}$ under a scheduled public
key. This object need not itself be an entry of $L$. The honest role supplies
an algorithm $A$ which derives $C_{pk}$, decrypts its local state with $sk$,
examines $L$, and returns $sk$ if the role is activated and $\bot$ otherwise.
If it returns $sk$, everyone can check $\mathsf{KeyMatch}$, decrypt
$C_{pk}$, and continue the public computation. Thus publishing the key is the
role's only message and models complete post-speech exposure. If $A$ returns
$\bot$, the secret key and encrypted role state remain hidden forever.

For this interpretation, every secret used by an ordinary role must be
recoverable from $sk$ and public protocol data. A protocol that gives the
role additional hidden auxiliary state cannot claim that publishing only
$sk$ models leakage of its complete state.

\paragraph{Strong input role.}
An application party $P$ embeds its private input in $A$ and causes the
ledger to append only the prescribed value $A(sk_P,L)$ and its signature.
For example, $A$ may decrypt two activation tokens and return the one selected
by an input bit. If $P$ remains honest, neither $sk_P$, the unused token, nor
the input-dependent algorithm is revealed. If $P$ is corrupted, the
adversary receives its ordinary UC state together with the key and protected
call history retained by the ledger; correspondingly the simulator receives
the input leakage specified by the target functionality. Corruption after
the turn does not permit a replacement message because $pk_P$ is already
spent.

The present finite construction has public output. A later private-output
variant may add a dual protected-read operation which evaluates a function of
$(sk_P,L)$ and returns only its result privately to $P$; no such operation is
needed here.

\subsection{Random corruption and finite committees}

The random coins apply only to computation-party registrations. Let $C$ be
any size-$k$ computation committee selected from registration indices without
using the hidden coins $(c_i)_i$, and let
$X=|\{i\in C:c_i=1\}|$. Then $X\sim\mathsf{Binomial}(k,p)$. At the concrete
choice $p=1/4$, Hoeffding's inequality gives
\[
  \Pr[X\geq k/3]\leq \exp(-k/72).
\]
If a finite execution uses $q$ such committees, a union bound gives
\[
  \Pr[\neg\mathsf{Good}]\leq q\exp(-k/72),
\]
where $\mathsf{Good}$ is the event that every committee has at most
$t=\lfloor(k-1)/3\rfloor$ corrupt members. Hence
$k=\Theta(\lambda+\log q)$ makes the bad event negligible for every
polynomial-size finite execution. Application parties do not enter this
calculation.

\subsection{Interpretation and limitations}

The preliminary experiment should not be read as saying that a real adversary
naturally corrupts machines by tossing independent coins. The intended
real-world picture has a large population of apparently interchangeable
machines and an adversary that may choose any fixed $p$ fraction of that
population to compromise. Fresh computation clients, keys, and role states
are placed on machines at random, without revealing the placement before a
machine is corrupted or speaks. A key's logical role may be public even though
the physical machine carrying its secret state is not linkable from the
ledger. Random placement therefore turns the adversary's fixed physical set
into an approximately random corrupt subset of each computation committee.
The independent Bernoulli coins in the present model collapse this placement
argument into the cleaner statement that each computation PID is corrupt
independently.

This is a fixed-population, non-mobile abstraction. Once a computation PID is
sampled honest it cannot later be corrupted. An adversary that can repeatedly
replace the corrupted quarter of the physical population is strictly stronger
and is not modelled here. Named application parties are different on purpose:
an account owner or other designated input provider is not plausibly hidden by
random placement, so it retains the ordinary adaptive corruption interface.
The random computation corruptions are an implementation assumption to be
simulated away, not a corruption rule for the target application.

The statement that the ledger stores $sk$ and evaluates $A(sk,L)$ is likewise
an ideal-interface device, not a proposal that a blockchain hold users' secret
keys or execute their private code. In the intended implementation, a machine
stores its own key and locally computes the one value it places on the public
board. Moving that local step behind the ledger interface prevents an honest
party's secret key from becoming an output visible to the UC environment. For
a computation role, publishing $sk$ is an explicit way to model that speaking
exposes all of its encrypted role state; for a named application party, only
the prescribed value is published unless the party is actually corrupted.

This abstraction is intentionally stronger and cleaner than such a system.
It assumes honest key generation, independent placement, no selective abort
of an unfavourable registration, instantaneous common knowledge of $L$,
unpreventable honest appends, exact synchronous timeouts, protected
evaluation under a secret key, and a protocol-dependent turn function fixed
at initialization. It does not model consensus, denial of registration, network
partitions, grinding, compromise correlations, or the mechanism that places
a computation key on a real machine. These assumptions must not be mistaken
for properties already obtained by a concrete blockchain.

Finally, the computation PIDs and their random corruption coins exist only in
the real or ledger-hybrid execution. The target ideal functionality has the
usual application interface: named input and output parties, their prescribed
adaptive corruptions, and the intended outputs. It contains no committees,
randomly corrupted computation population, garbled circuits, or
$\mathsf{WhosTurn}$. A security proof must simulate all of that protocol
machinery from the target functionality's leakage.

The final paper may keep only this interface and interpretation in the main
text. Once the API is stable, an appendix can formulate
$\mathcal F_{\mathsf{Ledger}}$ in UC notation and make precise ITM creation,
private delivery and retention of $A$, synchronous timeouts, and corruption
bookkeeping. Those details should express the API above rather than determine
it.

\section{The Basic Finite Construction}
\label{sec:basic-construction}

We illustrate the central idea for one evaluation of a public Boolean circuit
$F$.  The input providers are named application parties and the output is
public.  Computation roles are assigned to registered computation keys by a
public rule.  The construction uses the ledger only through the generic API
of \cref{sec:model}; garbling, role assignment, encryption, slot closure, and
message validation are protocol operations.  Our first proof target fixes the
corrupt set of named input providers before setup.  The adaptive interface of
\cref{sec:model} is retained as a later target.

This section is a bounded, two-phase blueprint.  It cleanly isolates the use
of silence, but does not by itself instantiate the sampler that correlates its
pre-input and post-input phases.  That requirement becomes property~$Z$ in
\cref{sec:proof-outline}; the streaming candidate of \cref{sec:ongoing}
implements the same schedule repeatedly.

\subsection{Public structural branches}

Fix a ledger prefix $L_0$ containing the registered input providers and at
least $2mk$ fresh computation keys, where $m$ is the number of input wires.
For each wire $w$, the public role map selects two disjoint size-$k$
committees $C^w_0$ and $C^w_1$.  All $2m$ committees are mutually disjoint,
and no computation key is assigned another canonical role.  Their subscripts
are \emph{structural branch indices}; they do not denote semantic bits.

For concreteness, the generic setup distribution may place a finite sampler
parameter $U_{\mathsf{fin}}$ in $\mathsf{crs}$.  Before inputs arrive, its
preparation phase samples for each input wire
\[
  (K^w_0,K^w_1),\qquad (T^w_0,T^w_1),\qquad
  \pi_w\sample S_2 .
\]
The garbling uses the semantic assignment
\[
  b\longmapsto K^w_{\pi_w(b)}.
  \label{eq:semantic-assignment}
\]
Equivalently, one may view it as garbling the structurally reindexed circuit
\[
  F_{\boldsymbol\pi}(z)=
  F(\pi_1^{-1}(z_1),\ldots,\pi_m^{-1}(z_m)).
\]
The pre-input public object $\mathsf{Prep}_F$ contains
the encrypted role packets described below, but not yet the garbled tables.
The complete label pairs and semantic maps remain hidden in the finite
sampler's logical state.

For public token validation, preparation publishes the branch-indexed tags
\[
  V^w_s=h_{\mathsf{tok}}(\mathsf{sid},w,s,T^w_s),
  \qquad s\in\{0,1\}.
  \label{eq:token-tags}
\]
Tokens have high min-entropy, and $h_{\mathsf{tok}}$ is assumed one way and
collision resistant on the domain-separated token inputs.  Thus the tags
allow everyone to recognize structural branch $s$ without revealing the
token in advance or the semantic value of that branch.

Label $K^w_s$ is verifiably $(t+1)$-out-of-$k$ shared among $C^w_s$, where
$t=\lfloor(k-1)/3\rfloor$.  For each member with public key $pk_i\in C^w_s$,
$\mathsf{Prep}_F$ contains
\[
  D_i=\mathsf{Enc}_{pk_i}
      (w,s,T^w_s,\mathsf{share}_{i,s},
       \mathsf{verification\ data}).
  \label{eq:finite-role-packet}
\]
It similarly contains, under the registered key of the application party
providing wire $w$, an encrypted packet specifying the map
$b\mapsto T^w_{\pi_w(b)}$.  Packet lengths, ordering, verification data, and
ciphertext distributions are symmetric in the two structural branches.
Neither $\pi_w$ nor the semantic bit represented by a public committee is
revealed.

The fixed $\mathsf{WhosTurn}$ waits until $L_0$ exists and then schedules the
provider keys and role keys in a canonical order.  Registration itself has no
role argument.

\subsection{Input activation and irrevocable default}

When the provider for wire $w$ is scheduled, an honest provider with input
$x_w$ supplies a protected algorithm $A_{w,x_w}$ to the ledger.  The
algorithm derives $\mathsf{Prep}_F$, decrypts the provider packet using
$sk_P$, and returns the tagged structural token $T^w_{\pi_w(x_w)}$.  The
ledger appends this output without revealing $sk_P$, the other token, or the
input-dependent algorithm.

A corrupt provider may instead put any string in its turn.  At a fixed cutoff
the protocol derives one irrevocable close record $m_w$.  Define
\[
  \mathsf{Valid}_w(m)=s
  \quad\Longleftrightarrow\quad
  h_{\mathsf{tok}}(\mathsf{sid},w,s,m)=V^w_s
\]
for exactly one $s\in\{0,1\}$; otherwise set
$\mathsf{Valid}_w(m)=\bot$.  A $\mathsf{NoMsg}$ record, malformed string,
ambiguous match, and late entry all give $\bot$.  The effective structural
branch is
\[
  s_w=\begin{cases}
       \mathsf{Valid}_w(m_w),&\mathsf{Valid}_w(m_w)\neq\bot,\\
       0,&\mathsf{Valid}_w(m_w)=\bot.
      \end{cases}
  \label{eq:finite-default}
\]
Because $\pi_w$ is uniform and hidden, default branch $0$ has semantic value
$\pi_w^{-1}(0)$, a uniform bit.  A corrupt provider that has opened its own
packet knows the semantic map and may choose its effective input, just as a
corrupt ideal input party may.

The roles, not a protocol-aware ledger, implement this rule.  Committee
$C^w_s$ speaks precisely if $\mathsf{Valid}_w(m_w)=s$, or if
$\mathsf{Valid}_w(m_w)=\bot$ and $s=0$.  In particular, $C^w_0$ can
distinguish a valid branch-$1$ token from a malformed token using the public
tags and therefore does not speak in the former case.  Closure must precede
the default release: otherwise a malformed message followed by a late valid
token could expose both labels.

\subsection{Conditional key release and evaluation}

Each computation member is scheduled once.  Its honest protected algorithm
derives its canonical $D_i$, decrypts it with $sk_i$, checks the public close
record and token tags, and returns $sk_i$ exactly under the rule above;
otherwise it returns $\bot$.  The selected committee therefore reveals its
keys, while every honest member of the sibling committee produces only a
$\mathsf{NoMsg}$ timeout and keeps its key hidden forever.

The ledger performs no validation.  Protocol observers ignore a purported
key unless $\mathsf{KeyMatch}(pk_i,sk_i)$ holds, decrypt $D_i$ themselves,
and reject malformed shares using the public verification data.  Conditioned
on $\mathsf{Good}$, the selected committee has enough honest members to reveal
at least $t+1$ valid shares despite corrupt withholding.  Everyone
reconstructs exactly one active structural label: $K^w_{\pi_w(x_w)}$ for a
valid honest input and $K^w_0$ for a defaulted input.

Only after all input slots close does the completion phase emit the public
garbled tables $\widetilde F$ and output-decoding information, using the
label pairs and assignment in \eqref{eq:semantic-assignment}.  All parties
evaluate $\widetilde F$ from the reconstructed active labels and decode the
public output.  This order is important for the intended straight-line
simulation: the ideal public output is available before the tables must be
fixed.  Publishing the tables before input closure would instead require an
online or adaptively simulatable garbling primitive.

Ordinary garbling decoding is sufficient.  A protocol needing a private bit
$y$ for an honest recipient can, at the application layer, let the recipient
input a fresh pad $r$ and publicly output $y\oplus r$; version~1 makes no
claim for later adaptive corruption of that pad holder.

\subsection{Why complete exposure need not be fatal}

Publishing $sk_i$ exposes the selected role's complete retained state,
including its packet plaintext, local coins, and call record.  The exposed
share belongs to a label reconstructed publicly anyway and gives no share of
the disjoint sibling label.  In the sibling committee, at most $t$ members are
corrupt on $\mathsf{Good}$, while every honest member remains silent;
threshold privacy therefore protects the unselected label.

The hidden permutation is necessary in this public-role model.  Without it,
the speaking committee would reveal the input.  With it, the public structural
branch has no fixed semantic meaning, and garbling privacy handles the
remaining dependence on the input.

The static-corruption restriction on named providers is used here.  A
provider corrupt before setup controls its slot; an honest provider is never
later asked to explain the encrypted unused token or its hidden semantic map.
An adaptive extension will require an equivocal or non-committing treatment of
the provider packet.

This remains a blueprint rather than a theorem.  A proof must combine the
two-stage finite-sampler property, multi-user encryption, token-tag security,
verifiable-sharing privacy and robustness, garbling privacy, and the
good-committee bound.  It must also account for retrospective decryption.
One-role-per-key gives one canonical packet, not one ciphertext globally:
adversarial evaluations of a copyable sampler may address further fork
packets to the same $pk_i$.  Canonical-root isolation must ensure that every
such plaintext is independent fork material containing no canonical inactive
label; the readout part of property~$Z$ must simulate all openings.

\paragraph{A gate-by-gate variant.}
For intuition, garbling can be replaced by an expensive
information-theoretic construction.  For every NAND gate and public
structural branch corresponding to a pair of input tokens, a separate
committee holds shares of the corresponding output token.  Only the committee
whose branch matches the public inputs reveals its keys.  This makes the role
of silence especially transparent, at the cost of a large setup and many
registered computation parties.

\section{A Streaming Garbled-Circuit Factory}
\label{sec:ongoing}

The finite construction samples one bounded garbling.  For ongoing
computation we now describe a more specific target than a general universal
streamer: an \emph{oblivious garbled-circuit sampler} (OGCS) for one fixed
public transition circuit
\[
  F(s_e,a_e)=(s_{e+1},y_e),
\]
where $s_e$ is carried state, $a_e$ is fresh input, and $y_e$ is public.  One
application of $F$ is a \emph{stretch}.  The desired public object has a finite
description and accepts arbitrarily long finite histories; in particular, its
interface has no setup-time horizon.

This section gives a candidate construction and a precise interface against
which a proof can be written.  It is not yet a theorem.  The cited
general-purpose ordinary-iO-based succinct TM-obfuscation results do not
immediately provide the horizon-free interface, and the final online
random-oracle programming step also remains open in this first version.

\subsection{The prefix interface and its schedule}

Setup publishes a finite parameter $U_{\mathsf{ogcs}}$ fixing $F$ and
containing a fresh identifier $\mathsf{uid}$.  The public algorithm is
evaluated on a complete finite input history:
\[
\begin{split}
  \mathsf{OGCS}(U_{\mathsf{ogcs}};I,
     &(M_0,\Sigma_0,\mathbf W_0),\ldots,\
       (M_e,\Sigma_e,\mathbf W_e))
       \longrightarrow D_e .
  \label{eq:ogcs-interface}
\end{split}
\]
The physical object is copyable and stateless; its hidden \emph{logical}
state is reconstructed by replaying the input history.  On every accepting
history its outputs are exactly prefix preserving,
\[
  D_{e-1}\preceq D_e,
\]
and repeated evaluation on the same history gives the same bit string.  A
failed check enters an absorbing reject state and emits no further extension.
The certificate is an acceptance witness, not stream input: two histories are
called command-equivalent when they have the same $I$ and the same sequence of
$(M_e,\mathbf W_e)$ values, even if they use different valid encodings of
$\Sigma_e$.  The program normalizes away certificate bytes after verification.
A different initialization or non-equivalent accepted command history is a
fork and must not expose the hidden state of the canonical history.

The initialization input $I$ produces a preparation object
$\mathsf{Prep}_0$ but no garbled table or application output.  For a fixed
public initial state $s_0$, it also publishes exactly one active initial-state
label per state wire as part of $\mathsf{Prep}_0$.  The first certificate key
$\mathsf{VK}_0$ is likewise a component of $\mathsf{Prep}_0$.  In general,
$\mathsf{Prep}_e$ is public before the input
slots of stretch $e$ open.  It contains
\begin{itemize}
  \item encrypted activation tokens and verifiable label shares, together
    with public token-validity tags, for the fresh external-input wires of
    stretch $e$; and
  \item encrypted branch secrets and public verification material for a
    fresh one-time certificate used to close that stretch.
\end{itemize}
The hidden plaintext material underlying $\mathsf{Prep}_e$ was fixed by the
preceding accepted extension, or by $I$ when $e=0$.  For notation set
$D_{-1}:=\mathsf{Prep}_0$.

The order within a stretch is important.
\begin{enumerate}
  \item Input providers post activation tokens.  At a deterministic ledger
    cut every slot is closed and the selected input committees reconstruct one
    active external-input label per wire.  Missing or malformed slots use the
    publicly detectable random-default rule of
    \cref{sec:basic-construction}.
  \item The registration window for stretch $e+1$ closes.  The canonical
    ledger cut determines one ordered vector $\mathsf{PKBatch}_{e+1}$.  If the
    required number of fresh keys is unavailable, the extension waits.
  \item Everyone computes the finite vector $\mathbf W_e$ of random-oracle
    values specified below.  This vector may be transported by any party and
    need not be stored on the ledger; it has no authority of its own.
  \item The certificate roles derive one canonical checkpoint $M_e$ from the
    closed ledger view and reveal a certificate $\Sigma_e$ for it.
  \item The accepted triple $(M_e,\Sigma_e,\mathbf W_e)$ appends
    \[
      \GC_e,\qquad \mathsf{Dec}_e,\qquad \mathsf{Prep}_{e+1}
      \label{eq:stretch-output}
    \]
    to the cumulative OGCS output.
  \item Everyone evaluates $\GC_e$.  This gives the public output $y_e$ and
    one active label on every next-state wire.
\end{enumerate}
Publishing the garbled tables only after input closure lets the straight-line
simulator obtain the public ideal output before it must fix those tables.  A
variant publishing the complete garbling earlier would require a stronger
online or adaptively simulatable garbling notion.

\subsection{A one-time common-input certificate from silence}
\label{sec:yamso-certificate}

The ledger has no signing key.  Persistence instead comes from fresh YAMSO
roles that certify a value they compute from their common ledger view.

Let
$\mathsf{Code}:\{0,1\}^{\kappa}\rightarrow\{0,1\}^{n}$ be a binary code of
relative distance $\delta$, let $h_{\mathsf{sig}}$ be collision resistant,
and let $f$ be one way.  For certificate epoch $e$, the hidden generator
samples independently
\[
  A^e_{j,b}\sample\{0,1\}^{\lambda},\qquad
  Y^e_{j,b}=f(\mathsf{sid},e,j,b,A^e_{j,b})
\]
for $(j,b)\in[n]\times\{0,1\}$.  The public verification material is
$\mathsf{VK}_e=(Y^e_{j,b})_{j,b}$.  Each preimage $A^e_{j,b}$ is encrypted to
a different fresh one-use computation role $P^e_{j,b}$.

At the fixed certification cut, every honest role independently computes the
same canonical message $M_e$ and the codeword
\[
  c_e=\mathsf{Code}(h_{\mathsf{sig}}(M_e)).
\]
The protected algorithm of role $P^e_{j,b}$ returns its secret key precisely
when $c_e[j]=b$.  Observers check $\mathsf{KeyMatch}$, decrypt the canonical
role packet, and treat the recovered $A^e_{j,b}$ as the opening.  A
certificate contains valid openings in at least $\tau n$
coordinates matching the claimed codeword.  Corrupt matching roles may
withhold, which is why verification does not require every coordinate.

Set $\rho=p$, the computation-role corruption probability of the model, and
choose constant margins
$\gamma_{\mathsf{live}},\gamma_{\mathsf{forge}}>0$ such that
\[
  1-\delta(1-\rho)+\gamma_{\mathsf{forge}}
  <\tau<
  1-\rho-\gamma_{\mathsf{live}}.
  \label{eq:certificate-window}
\]
The right inequality gives liveness.  For a different codeword, every
disagreeing coordinate needs the opposite opening, which is available without
inverting $f$ only when that opposite role is corrupt; the left inequality
prevents enough such openings.  The interval is possible when
$\delta>\rho/(1-\rho)$, and hence permits any $\delta>1/3$ at
$\rho=1/4$.  The canonical digest may depend on corrupt-role information
visible before its cut.  We therefore require the good-assignment event to
hold simultaneously for all possible canonical codewords and all ordered
pairs of distinct codewords.  A conservative sufficient condition is
\[
  2^{\kappa}\exp(-\Omega(\gamma_{\mathsf{live}}^2 n))
  +2^{2\kappa}\exp(-\Omega(
       \gamma_{\mathsf{forge}}^2\delta n))
\]
to be negligible.  The exact constants and concentration proof are part of
the certificate lemma; they are not claimed proved here.

Speaking exposes the complete state of $P^e_{j,c_e[j]}$, but its branch secret
has just become public.  The role containing $A^e_{j,1-c_e[j]}$ remains
silent forever.  This is precisely where the at-most-once rule creates
persistent authenticity without a persistent signing key.

The distance-amplified Lamport idea is not new.  Katz and Vaikuntanathan use a
high-distance encoding in a leakage-resilient one-time signature
\cite{KatzVaikuntanathan2009LeakageSignatures}.  We reuse the distance
intuition, while the distributed role assignment, partial liveness threshold,
post-speech full exposure, and common-input property require a separate
argument.  Kelsey, Lang, and Lucks study distributed stateful hash-based
signing and the disjoint-honest-trustees one-time-key-reuse problem in the
absence of broadcast \cite{KelseyLangLucks2025DistributedHBS}.

The certificate is not a signing service for adversarially supplied messages.
Honest roles derive $M_e$ from one already closed ledger cut.  The common
ledger prevents them from being shown different inputs; the certificate makes
that public truth usable by an obfuscated program that cannot read the ledger
itself.

\subsection{Keyless bootstrap and certified commands}

The initialization string has the form
\[
  I=(\mathsf{Init},\mathsf{uid},\mathsf{sid},
       \mathsf{initDigest},\mathsf{PKBatch}_0).
\]
The program requires the supplied $\mathsf{uid}$ to equal the one embedded in
$U_{\mathsf{ogcs}}$.  The canonical $\mathsf{initDigest}$ commits to the
initial ledger cut, public initial state, role allocation, and all relevant
parameters.  The hidden program absorbs the complete encoding of $I$ before
deriving its root and emitting $\mathsf{Prep}_0$, which includes
$\mathsf{VK}_0$ and the active labels for the fixed $s_0$.  Thus initialization
is not literally outputless: it emits branch-local encrypted preparation
material and public labels for a public state, but no garbled table or
application-derived output.

Anyone may run the copyable program on $I'\neq I$.  Such a call may even use
only adversarial public keys and hence create an adversarially controlled
certificate chain.  This is harmless only because the different
initialization is absorbed before any branch material is derived: every
token, label, ciphertext, and later certificate key on that fork lies under a
pseudorandomly independent logical root.  Later exposure of a canonical role
key may decrypt a packet addressed to the same public key on such a fork, but
the plaintext contains only fork material.  The common ledger selects the
live root; no certifying-ledger secret key is needed.

For a public partition of microstep indices into finite intervals $J_e$, let
\[
  W_i=H(\mathsf{ogcs},\mathsf{uid},\mathsf{sid},i),
  \qquad i\in J_e,
\]
write $\mathbf W_e=(W_i)_{i\in J_e}$, and set
$\omega_e=\mathsf{Hash}(J_e,\mathbf W_e)$.  The obfuscated program has no
random-oracle access, so these values are computed outside it.  Honest
certificate roles compute the same vector themselves.  The vector supplied
with the OGCS call is nonauthoritative: the program checks its digest against
the certified $\omega_e$.

Define the full pre-oracle descriptor and its digest by
\[
  Q_e=(\mathsf{uid},\mathsf{sid},e,
    \mathsf{Hash}(D_{e-1}),\mathsf{Hash}(\mathsf{Prep}_e),
    \mathsf{inputCheckpoint}_e,\mathsf{PKBatch}_{e+1}),
  \qquad
  \chi_e=\mathsf{Hash}(Q_e),
\]
and let the canonical certificate message be
\[
  M_e=(Q_e,J_e,\omega_e).
  \label{eq:checkpoint}
\]
The next-key registration cutoff occurs before this message is certified.
Exact closed-ledger positions, role allocation, and all domain separators are
part of $Q_e$'s canonical encoding.  Thus the program receives, checks, and
can use the actual input checkpoint and recipient-key batch rather than only
their hash image.  It checks the predecessor, epoch, syntax, oracle-vector
digest, and certificate under the verification material in its current
logical frontier.  Any failure enters the absorbing reject state.

The next verification key is an output inside $\mathsf{Prep}_{e+1}$ and hence
is authenticated by the current accepted extension.  It is not included as an
independent input to the same command, avoiding a circularity.  Different
valid subsets or orderings of openings for the same $M_e$ induce the same
extension: all derivations use $M_e$, never the certificate bytes.

\subsection{The hidden logical stream}

Let $\mathsf{InMat}_e$ contain the external-input label pairs, activation
tokens, semantic permutations, shares, and role-packet plaintexts underlying
$\mathsf{Prep}_e$, and let $\mathsf{CertMat}_e$ contain its certificate
preimages and packet state.  At the beginning of stretch $e$, the logical
frontier is
\[
  \Phi_e=\left(
    X_e,
    \mathsf{StatePairs}_e,
    \mathsf{InMat}_e,
    \mathsf{CertMat}_e,
    \mathsf{VK}_e
  \right),
  \label{eq:logical-frontier}
\]
After the public cuts are fixed, let $\overline\Phi_e$ be the frontier
$\Phi_e$ with its seed coordinate replaced by
\[
  \overline X_e=
    \mathsf{KDF}(X_e,(\mathsf{publicCut},\chi_e)).
\]
Write $\Phi_{e,i}$ for the corresponding intermediate frontier after
microstep $i$, with $\Phi_{e,\ell_e-1}=\overline\Phi_e$ and
$\Phi_{e,r_e}=\Phi_{e+1}$ on a normally generated stretch.  Here
$\mathsf{StatePairs}_e=
 \{(K^{\mathsf{st}}_{e,w,0},K^{\mathsf{st}}_{e,w,1})\}_w$.
The active member of every state pair is public from evaluation of the
preceding stretch.  For $e=0$, $\mathsf{Prep}_0$ publishes the member encoding
the fixed public initial state; the sibling remains hidden.

Enumerate the finite objects generated for the next suffix by the consecutive
indices $J_e=[\ell_e,r_e]$.  Set
\[
  X_{e,\ell_e-1}=\overline X_e,
  \qquad
  X_{e,i}=\mathsf{KDF}(X_{e,i-1},W_i)
  \quad(i\in J_e),
  \qquad
  X_{e+1}=X_{e,r_e}.
  \label{eq:kdf-chain}
\]
These equations specify the normal branch.  At a programmed microstep, a
valid capsule replaces both the emitted fragment and the complete successor
microfrontier $\Phi_{e,i}$; subsequent generation continues from that
replacement frontier.
Domain-separated puncturable-PRF evaluations from these states derive the
output-state label pairs and garbling coins for $\GC_e$, followed by
$\mathsf{InMat}_{e+1}$, $\mathsf{CertMat}_{e+1}$, verification material,
verifiable shares, and role-state encryption coins for
$\mathsf{Prep}_{e+1}$.  The current $\mathsf{InMat}_e$ is not regenerated:
it is the material already committed before the input window.  The complete
pre-oracle cut is therefore absorbed before the first microstep.

Readout for the newly prepared input roles is bound to the exact ordered
recipient-key vector.  For example, its coins have the form
\[
 r^{\mathsf{read}}_{e+1,w,s,j}
 =\mathsf{PPRF}_{K_{\mathsf{read}}}
   (X_{e,i},\chi_e,e+1,w,s,j,
    pk_{e+1,w,s,j}).
  \label{eq:pk-bound-readout}
\]
Consequently, a canonical packet is tied to its intended recipient key.
Packets produced by reset or fake-initialization calls under the same public
key are not claimed never to decrypt; rather, prefix isolation must ensure
that their plaintexts contain only independently rooted fork material.

The garbling $\GC_e$ uses $\mathsf{StatePairs}_e$ and the external-input pairs
in $\mathsf{InMat}_e$.  Its output-state labels are literally the pairs in
$\mathsf{StatePairs}_{e+1}$ stored in $\Phi_{e+1}$.  Evaluation makes only the
active member of each pair public.  No machine links an old label to a new
one, and no committee is needed for carried state.

Changing $I$, an accepted certified checkpoint, or any earlier $W_i$ changes
the logical derivation for the suffix after that point.  On the live protocol
path, the certificate prevents such a change after a shared state label
exists.  Calls on another initialization root remain independently evaluable
but receive no honest input material belonging to the live root.  In a
programmed proof hybrid, a sequence of transition capsules may prescribe each
affected public fragment together with its matching successor microfrontier;
no capsule may program only a table while leaving an inconsistent hidden
successor.

\subsection{Candidate realization by delayed programming}

For a declared input-length bound $L$, a finite version of the public object
may be based on succinct iO for bounded-input Turing machines.  Cryptographic
iterators and succinct TM-iO provide relevant techniques for compactly
representing long bounded computations
\cite{KoppulaLewkoWaters2014IOTM,AnanthJainSahai2017TMIO}.  They do not by
themselves instantiate the horizon-free interface above, whose replay input
has unbounded length; that requires an additional primitive or a
reformulation of the interface.  We leave this as a separate realization gap
rather than hide $L$ in the environment's running-time polynomial.

For a bounded prefix, we propose to adapt the delayed-backdoor programming of
Hofheinz et al.'s universal sampler
\cite{HofheinzEtAl2014UniversalSamplers}.  Its hidden-trigger paradigm, and
the later selectively secure pseudorandom puncturable deterministic encryption
(PPDE) of Koppula, Poelstra, and Waters
\cite{KoppulaPoelstraWaters2017FastSamplers}, let an apparently random string
carry a programmed bounded object.  In our proposed adaptation, a random
$W_i$ selects the ordinary PRF-generated transition.  In a proof, the oracle
answer at one punctured point is changed to encode a prescribed transition
capsule.

More explicitly, let ($\mathsf{BEnc},\mathsf{BDec}$) denote the desired
puncturable carrier.  Its ciphertexts are pseudorandom, while a uniform string
decrypts to $\bot$ except with negligible probability.  The same hidden
microstep program is used in both the real execution and the proof:
\[
  C_i\leftarrow\mathsf{BDec}_{K_{\mathsf{bd}}}(W_i),
  \qquad
  \mathsf{Step}_i=
  \begin{cases}
    C_i,&\text{if }\mathsf{ValidCap}(C_i,\Phi_{e,i-1},\chi_e)=1,\\
    \mathsf{NormalGen}(X_{e,i-1},W_i,\chi_e,i),&\text{otherwise.}
  \end{cases}
  \label{eq:normal-programmed-step}
\]
On honest random-oracle answers, the normal branch is taken except with
negligible probability.  In a selective bounded-prefix hybrid in which
$C_i$ is fixed before carrier-key setup, the intended proof first replaces
the embedded carrier key by a key punctured at $C_i$ and hardcodes the
prescribed behavior at the exceptional ciphertext.  It separately punctures
the normal-generation keys at the same logical point.  Koppula--Poelstra--
Waters PPDE pseudorandomness can then hide the swap from uniform $W_i$ to
$\mathsf{BEnc}_{K_{\mathsf{bd}}}(C_i)$.  Choosing $C_i$ online after
carrier-key setup requires an adaptive carrier argument and remains open.

The programmed branch bypasses the normal KDF and readout derivations for
that microstep and takes the complete successor microfrontier from its
capsule.  Thus, in the selective experiment, the capsules can be sampled from
$\chi_e$, the closed recipient keys, and the prescribed output before $W_i$
or $M_e$ is fixed.  This is a specification of the carrier property we need,
not a claim that the cited PPDE theorem already supplies the complete stateful
interface.

For our use each capsule must contain the \emph{complete} affected logical
microtransition: its public garbled fragment and every hidden successor value
on which later fragments depend.  Programming only one public gate while
leaving an inconsistent successor $X_{e,i}$ would not preserve the stream.
The intended hybrid punctures the KDF/PRF path at one canonical microstep,
installs this complete capsule, leaves every sibling subtree pseudorandom,
and then repeats from left to right.

Two further proof obligations remain visible.  First, computational
unforgeability of the YAMSO certificate does not make two obfuscated programs
functionally identical on every mathematical input.  In their binding setup,
Fernando et al.'s somewhere statistically unforgeable signatures make
alternate valid signatures at a selected round and message statistically
nonexistent, except with negligible probability over setup.  They use this
property to establish functional equivalence of adjacent obfuscated gate
programs before applying ordinary iO \cite{FernandoEtAl2023SSU}.  This is a
candidate technique here, not an instantiation we may assume for free.

Second, the points
$H(\mathsf{ogcs},\mathsf{uid},\mathsf{sid},i)$ are predictable after
$U_{\mathsf{ogcs}}$ is public.  A selective-capsule experiment can fix the
target $\mathsf{sid}$, finite index set $\bigcup_eJ_e$, canonical schedule,
and complete capsule sequence before publication, program all corresponding
points, and then rely on the fresh $\mathsf{uid}$ to rule out an earlier
query.  This is structural evidence for the bounded carrier hybrid, not a
real--ideal simulation: in the protocol the honest inputs and outputs are not
known before setup.  It also does not give one online simulator for the
horizon-free execution.

The intended next pass prepares a fresh unpredictable nonce one stretch
ahead.  In that proof the simulator first waits for the input and next-key
cuts, receives $y_e$ privately, and uses $Z.\mathsf{Complete}$ and
$Z.\mathsf{Prepare}$ to form the microstep capsules from $\chi_e$, the
recipient keys, and $y_e$.  It then programs the nonce-dependent oracle
points, reveals and recomputes $\mathbf W_e$, forms $M_e$, and only then
releases the honest openings that take $\Sigma_e$ across its acceptance
threshold.  Thus no one can invoke the accepted canonical extension before
all of its capsules are installed.
We do not repair either issue by giving the environment an explicit time bound.

\subsection{What the candidate establishes and what it does not}

The construction isolates a plausible path to indefinitely extensible
garbling.  Adjacent stretches share boundary labels by construction; new
public keys enter only through a certified cutoff; external oracle values are
bound by a certified digest; and the only online secret holders that speak
reveal branch-local state already made public.  A fake initialization can
generate a whole fake computation, but it forks before any live input or
state material.

The missing statements are now explicit: the one-time certificate lemma, a
two-stage finite garbling/readout simulation with prescribed boundary labels
and full speaker exposure, a horizon-free public-program realization, the
stateful delayed-programming lemma, exact iO-compatible authentication, and
late online oracle programming.  These are organized as whole-process
replacement lemmas in \cref{sec:proof-outline}.  A general universal-streamer
abstraction can be revisited after this fixed-function OGCS is understood;
the present construction does not claim one.

\section{A Modular Proof Outline}
\label{sec:proof-outline}

The final security statement should be phrased at the application interface.
We do not, however, decompose the intermediate argument into a collection of
UC ideal functionalities for garbling, OGCS, certificate preparation, and
committee readout.  A naive decomposition whose interfaces output the secret
keys and label pairs needed by the next hybrid would also expose those values
to the UC environment.  This is the proof-engineering problem encountered
when secret key material is passed through intermediate hybrids in the
asynchronous YOSO proof \cite{DamgardEtAl2025AsyncYOSO}.

We instead compare two complete interactive processes.  They include the
environment, adversary, random oracle, ledger, all evaluations of the copyable
public OGCS object, every protocol message, and the complete retained state of
every role after it speaks.  The replacement lemmas below are ordinary
indistinguishability statements about these whole executions.  A resulting
whole-process argument is intended to be used inside a separately constructed
top-level UC proof; such a wrapper is not automatic and is not claimed here.

\subsection{Scope of the first statement}

Before setup, the adversary chooses the corrupt set
$\mathcal C_{\mathsf{in}}$ of named input providers.  This set remains fixed.
Each computation registration $i$ is independently marked corrupt by
$c_i\sample\mathsf{Bernoulli}(p)$, and an honest computation role cannot later
be corrupted.  Nevertheless, immediately after an honest role speaks, its
complete role state is given to the adversary.  Every computation public key
has one canonical role, all outputs are public, and a closed missing or
malformed input has the random-default semantics of
\cref{sec:basic-construction}.

The desired OGCS construction accepts every finite history without a
setup-declared horizon.  The proof template below therefore assumes a
horizon-free online property~$Y$.  The bounded-input TM-iO and selective
oracle-programming evidence in \cref{sec:ongoing} establishes neither that
property nor a hidden running-time bound; both limitations remain explicit.

\subsection{The real and ideal processes}

The process
$\mathsf{Real}_{\Pi,\mathcal A,\mathcal Z}(1^\lambda)$ runs the protocol of
\cref{sec:model,sec:basic-construction,sec:ongoing}.  In addition to the
ordinary ledger view, $\mathcal A$ and $\mathcal Z$ receive
\begin{itemize}
  \item the complete initial state of every corrupt registered role;
  \item every public derived garbling, ciphertext, share opening, validity
    tag, and certificate;
  \item the result of every adversarial reset or fork evaluation of
    $U_{\mathsf{ogcs}}$; and
  \item the complete secret key, canonical packet plaintext, coins, and
    retained call record of every honest computation role immediately after
    its message is irrevocable.
\end{itemize}
The same disclosed key may retrospectively decrypt packets addressed to it by
off-path evaluations of the copyable object; those plaintexts are also part
of the real view.  Silent honest roles expose nothing and are never reused.

The process
$\mathsf{Ideal}_{F,\mathcal S,\mathcal Z}(1^\lambda)$ evaluates the intended
stateful transition directly.  Honest named providers give their inputs to
the ideal computation; in the first statement every honest provider supplies
exactly one bit before its slot is scheduled.  For a corrupt provider,
$\mathcal S$ supplies the semantic bit extracted from its uniquely closed
protocol slot, including a bit of the simulator's choice for a missing or
malformed slot.  It computes
\[
  (s_{e+1},y_e)=F(s_e,a_e)
\]
and, once both the input and next-key cuts are fixed, sends
$(\mathsf{OutputReady},e,y_e)$ privately to $\mathcal S$.  The simulator may
use this value internally while it prepares the certified extension, but does
not deliver it to $\mathcal A$ or $\mathcal Z$.  The ideal process releases
$y_e$ publicly only at the corresponding simulated evaluation event.  Thus
the environment sees the output at the same point as in the real execution,
while $Z.\mathsf{Complete}$ learns it causally.  The ideal process contains no
computation PIDs, committees, labels, garbled circuits, random corruption
coins, or OGCS state.  In the simulation, a default is coupled to
$\pi_w^{-1}(0)$; since $\pi_w$ is uniform, this has exactly the ideal default
distribution even though the ideal corrupt-party interface lets the simulator
supply the bit.

\subsection{Three replacement lemmas}

The proof is organized around three properties, denoted $X$, $Y$, and $Z$.
They are lemmas about the surrounding execution, not callable ideal
functionalities.

\paragraph{$X$: canonical certification from one-use exposure.}
Condition on the good certificate-role event.  For every epoch, the canonical
checkpoint obtains enough valid openings, while no PPT context obtains an
accepted certificate for a different checkpoint under the same verification
material.  This remains true after every role that certified the canonical
checkpoint exposes its complete state.

Property~$X$ is a non-equivocation-and-liveness statement.  Because the
canonical digest may depend on corrupt-role information visible before its
cut, the good-assignment event is required to hold uniformly for every
possible canonical codeword and every ordered pair of distinct codewords.
The simulation of encrypted certificate packets, role keys, coins, and
exposed call records is
assigned to $Z.\mathsf{Prepare}$ below.  The authentication proof separates a
collision in $h_{\mathsf{sig}}$, inversion of an unopened one-way image, and
an unusually favourable random corrupt-role pattern.  With constant margins
as in \eqref{eq:certificate-window}, a conservative rate-versus-tail
condition is
\[
  2^{\kappa}\exp(-\Omega(\gamma_{\mathsf{live}}^2 n))
  +2^{2\kappa}\exp(-\Omega(
      \gamma_{\mathsf{forge}}^2\delta n))=\negl(\lambda).
\]
Sequential composition then pays a union bound over the polynomial number of
certificate epochs appearing in the inspected execution.

Status: this lemma is \emph{conjectured}.  Its high-distance Lamport core is
source-backed by \cite{KatzVaikuntanathan2009LeakageSignatures}; the partial
opening threshold, common-input restriction, random role assignment, and full
post-speech exposure experiment still need a complete proof.

\paragraph{$Y$: programmable prefix-preserving OGCS.}
A horizon-free online simulator first produces one public OGCS parameter,
before the future initialization, certified descriptors, or outputs are
known.  When the canonical initialization $I$ arrives, it runs the normal
initialization derivation honestly, publishes $D_{-1}=\mathsf{Prep}_0$, and
retains the resulting hidden frontier $\Phi_0$.  This step requires no
programmed capsule because $s_0$ is public and there is no application-derived
output.  Once the
pre-oracle descriptor $Q_e$, the closed $\mathsf{PKBatch}_{e+1}$, and the
target microstep capsules for stretch $e$ are available, but before
$\mathbf W_e$ is exposed, it supports the simulation-side operation
\[
  \mathsf{Program}(Q_e,\mathsf{Capsules}_e)
      \longrightarrow\mathbf W_e.
\]
The canonical message
$M_e=(Q_e,J_e,\mathsf{Hash}(J_e,\mathbf W_e))$ is formed and certified only
afterwards.  The resulting object must satisfy
\begin{enumerate}
  \item every canonical call returns the prescribed output prefix;
  \item every accepting history is bit-for-bit prefix consistent, and
    repeated calls to one history agree exactly;
  \item each prescribed capsule contains the complete hidden successor
    microfrontier as well as its public output fragment;
  \item after the first non-equivalent divergence in $I$, $Q_e$, $J_e$, or
    $\mathbf W_e$, arbitrary fork and reset calls reject or evolve under
    pseudorandomly independent roots; valid certificate-witness variants for
    the same command are normalized away; and
  \item later disclosure of a canonical secret key may decrypt off-path
    packets addressed to that key, but all such plaintexts are simulatable
    independent-fork material containing no canonical inactive label or
    master seed.
\end{enumerate}
The surrounding context may interleave public OGCS evaluations with all
ledger activity and speaker exposures.

The candidate bounded-prefix reduction punctures the KDF/PRF path at the next
canonical microstep, uses delayed backdoor programming to install the capsule,
and replaces sibling subtrees by independent seeds.  A separate
authentication argument must establish exact or statistical functional
equivalence for each iO change; computational certificate unforgeability alone
is insufficient.

Status: this lemma is \emph{open}.  The cited general-purpose
ordinary-iO-based succinct TM-iO gives only a declared bounded-input object.
With predictable oracle points, the PPDE evidence applies only to a
selective-capsule bounded-prefix experiment in which the capsules are fixed
before carrier-key setup, not to the real--ideal simulation above.  The
hidden-nonce lift at the end of \cref{sec:ongoing} is the candidate online
repair for the latter gaps.

\paragraph{$Z$: two-stage finite preparation, garbling, and exposure.}
Property~$Z$ has an explicit causal interface.
\begin{description}
  \item[$Z.\mathsf{Prepare}$.]
    Before the inputs of stretch $e$ are known, fix structural external-input
    label pairs, token pairs, semantic permutations, public token-validity
    tags, corrupt shares, honest encrypted role packets, and the complete
    certificate preparation for that stretch.  For $e=0$, it consumes the
    honestly derived material returned by $Y$'s initialization step rather
    than resampling an incompatible frontier; that material includes full
    initial-state pairs and the published active member representing $s_0$.
  \item[$Z.\mathsf{Complete}(y_e)$.]
    After the input slots close, take the already fixed active external-input
    labels, the active carried-state labels, and the public ideal output
    $y_e$.  Produce garbled tables and decoding information consistent with
    those labels, together with fresh hidden next-state pairs and their active
    members.  The procedure need not learn the semantic input or carried
    state.  Its fresh preparation call for stretch $e+1$ uses the already
    closed $\mathsf{PKBatch}_{e+1}$ and pre-oracle descriptor $Q_e$.  This is
    a simulator-side target sampled without $\mathbf W_e$ or $M_e$;
    property~$Y$ must subsequently choose or program $\mathbf W_e$ so that the
    public object realizes that target.
\end{description}

Condition on every ordinary readout committee having at most $t$ corrupt
members.  Initially corrupt members' ciphertexts and shares must be fixed
honestly from preparation.  When a selected honest role later speaks, its
key, packet, token, coins, shares, and retained call record must be exposed
consistently with the reconstructed active label.  Encryption hybrids may be
used only under honest keys that remain silent.  One-role-per-key guarantees
one \emph{canonical} packet, not one ciphertext globally: packets created on
adversarial fork histories are handled by the isolation clause of $Y$, while
$Z$ simulates their later openings as fork material.  Threshold privacy hides
the canonical inactive label, and the hidden semantic permutation hides the
meaning of the observed structural branch.

The garbling sublemma must support prescribed active external and carried
labels, a prescribed public output, and a pair of hidden next-state labels
whose active member can be carried into the next simulated stretch.  This is
stronger syntax than a bare statement of ordinary one-shot garbling privacy.
Publishing the tables only after input closure is what makes
$Z.\mathsf{Complete}(y_e)$ causal.

Status: this lemma is \emph{conjectured}.  Static input-provider corruption is
used critically: an honest provider's unopened token map is never exposed
after the simulator has selected a dummy structural branch.  Token-tag
collisions or inversions, sharing failures, and malformed-opening acceptance
are recorded separately as auxiliary bad events.

\subsection{The straight-line simulator}

Assuming $X$, $Y$, and $Z$, the global simulator proceeds as follows.
\begin{enumerate}
  \item It samples registration and every immutable computation-corruption
    coin exactly as in the real process and starts the online simulator of
    $Y$.  When canonical $I$ is fixed, it invokes $Y.\mathsf{Initialize}(I)$
    to derive $\mathsf{Prep}_0$ and $\Phi_0$ honestly, then passes the
    corresponding input and certificate material to $Z.\mathsf{Prepare}$ for
    stretch $0$.  No initialization capsule is programmed.
  \item For a statically corrupt provider it forms the provider packet and
    token map honestly.  It extracts the effective input from the closed slot.
    On a malformed or missing slot it supplies $\pi_w^{-1}(0)$ to the ideal
    process, coupling the uniform protocol and ideal defaults.
  \item For an honest provider it publishes one uniformly selected valid
    structural token.  This distribution is independent of the provider's
    semantic bit, and static corruption means the unused semantic map is never
    opened.  The honest input reaches the ideal process without being shown to
    the simulator.
  \item It fixes corrupt committee shares from preparation.  For the selected
    honest committee it later exposes role states reconstructing the active
    label; honest sibling roles remain silent.
  \item Once the input slots and $\mathsf{PKBatch}_{e+1}$ have both closed,
    it forms the full pre-oracle descriptor $Q_e$, supplies all extracted
    corrupt inputs to the ideal process, and privately receives
    $(\mathsf{OutputReady},e,y_e)$.
  \item It invokes $Z.\mathsf{Complete}(y_e)$ with the active external and
    carried labels, and invokes $Z.\mathsf{Prepare}$ for stretch $e+1$ under
    the closed next-key batch.  These outputs determine a finite sequence of
    complete microstep capsules.  Collectively, their public fragments are
    $\GC_e$, decoding data, and $\mathsf{Prep}_{e+1}$; their final successor is
    the entire frontier for stretch $e+1$.  Before an accepted certificate
    exists, it calls
    $Y.\mathsf{Program}(Q_e,\mathsf{Capsules}_e)$ to install all of them and
    obtain $\mathbf W_e$.
  \item Only after all capsules for the extension are installed does it expose
    $\mathbf W_e$,
    form $M_e$, and use $X$ to produce the canonical certificate transcript.
    In particular, the final honest opening that crosses the acceptance
    threshold is released after installation.  Property~$X$ confines the
    extension to the unique certified checkpoint.
  \item Anyone may now evaluate the canonical OGCS extension.  Public garbled
    evaluation exposes the active next-state labels and triggers public
    delivery of $y_e$, after which the simulator repeats.
\end{enumerate}
At no point does the simulator receive the ideal hidden state or an honest
input bit.  It uses only corrupt inputs, prescribed post-speech leakage, and
public outputs.

\subsection{First bad event rather than circular conditioning}

Committee-overflow events are handled first by conditioning on independent
$\mathsf{GoodCom}$ and $\mathsf{GoodCert}$ events and paying their explicit
probabilities.  Within those conditioned experiments, let
$\mathsf{BadCrypto}$ be the first occurrence of a token-tag or digest
collision, inversion of an unopened one-way value, acceptance of a malformed
key or share, a false hidden trigger, or a premature oracle query in the
bounded selective programming experiment.  Protocol invariants---unique
canonical key assignment, absorbing rejection, and exact prefix
consistency---are requirements of the construction or property~$Y$, not
probabilistic adversarial events.

For each cryptographic bad event, we stop a hybrid immediately before its
first occurrence.  Adjacent stopped processes are shown indistinguishable;
the event is proved negligible in the endpoint where this is easiest; and
negligibility is transported back through the hybrids.  This is the
``Everywhere from Anywhere'' pattern of
\cite[Lemma~8]{DamgardEtAl2025AsyncYOSO}.  The transport loss is retained in
the bookkeeping rather than silently discarded.

\subsection{Conditional bookkeeping theorem}

Let $Q_{\mathsf{com}}$ be the number of ordinary size-$k$ readout committees
in the inspected prefix.  For $p<1/3$ and
$t=\lfloor(k-1)/3\rfloor$,
\[
  \delta_{\mathsf{com}}
  \leq Q_{\mathsf{com}}
       \exp\bigl(-2k(1/3-p)^2\bigr),
\]
which is $Q_{\mathsf{com}}e^{-k/72}$ at $p=1/4$.  Let
$\delta_{\mathsf{cert}}$ denote the distinct certificate-role concentration
and code-union-bound failure from $X$.  Let
$\epsilon_X,\epsilon_Y,\epsilon_Z$ be the three replacement advantages, and
let $\epsilon_{\mathsf{aux}}$ collect only disjoint collision, token-hash,
sharing, malformed-opening, and primitive-correctness errors not already
included in those terms.

\paragraph{Conditional whole-execution emulation template.}
If properties $X$, $Y$, and $Z$ hold with the horizon-free online quantifiers
stated above, then
\[
  \forall\mathcal A\;\exists\mathcal S\;\forall\mathcal Z:
  \qquad
  \mathsf{Real}_{\Pi,\mathcal A,\mathcal Z}
  \approx
  \mathsf{Ideal}_{F,\mathcal S,\mathcal Z}.
\]
A conservative transported-hybrid accounting has the form
\[
\begin{split}
  \mathsf{Adv}^{\mathsf{real\text{-}ideal}}_{\Pi,F}
  \leq{}& \delta_{\mathsf{com}}+\delta_{\mathsf{cert}}\\
        &+2\bigl(\epsilon_X+\epsilon_Y+\epsilon_Z
                 +\epsilon_{\mathsf{aux}}\bigr).
\end{split}
\]
The factor two is only a safe placeholder for transporting bad-event bounds;
a completed hybrid sequence should replace it by exact accounting.

This is deliberately a theorem \emph{template}, not a claimed security
theorem.  Property~$Y$ is open in both its horizon-free realization and online
oracle timing, property~$Z$ still requires a concrete two-stage
prescribed-label garbling lemma, and property~$X$ needs its distributed
full-exposure proof.  Writing the statement now fixes the interfaces between
the arguments without pretending that those arguments have already been
completed.


\bibliographystyle{splncs04}
\bibliography{bibliography,ogcs-references}

\end{document}
